\documentclass[11pt,a4paper]{article}
\usepackage[utf8]{inputenc}
\usepackage[T1]{fontenc}
\usepackage{graphicx}
\usepackage{geometry}
\usepackage{float}
\usepackage{hyperref}

\geometry{margin=2.5cm}
\setlength{\parskip}{0.7em}
\setlength{\parindent}{0pt}

\title{ Laporan Analisis Dataset: rendo\_electricity\_2013\_3 }
\author{ Local Agentic Analytics }
\date{\today}

\begin{document}
\maketitle

\begin{abstract}
Laporan ini menyajikan analisis dataset berdasarkan grafik agregat dan insight otomatis. Ringkasan temuan utama: Grafik distribusi net\_manager menunjukkan nilai yang konsisten dan seragam dengan rata-rata sebesar 8716912000008.0 kW. Analisis data energi menunjukkan rata-rata global active power sebesar 8716912000008 kW.
\end{abstract}

\section{Introduction}
Laporan ini menganalisis dataset 'rendo\_electricity\_2013\_3' yang disimpan pada tabel 'rendo\_electricity\_2013\_3'. Dataset memiliki 7 kolom numerik yang menjadi dasar visualisasi dan analisis statistik ringkas.

\section{Methodology}
Data terstruktur disimpan dan diagregasi menggunakan DuckDB agar analisis tetap ringan di lingkungan lokal. Visualisasi dibuat secara deterministik dengan matplotlib dari hasil query agregat, sedangkan narasi tiap grafik dihasilkan oleh InsightAgent berbasis model lokal menggunakan statistik ringkas, bukan data mentah.

\section{Detailed Analysis}
\subsection{ Distribusi net\_manager }
Grafik distribusi net\_manager menunjukkan nilai yang konsisten dan seragam dengan rata-rata sebesar 8716912000008.0 kW. Nilai minimum dan maksimum yang sama, mengindikasikan bahwa distribusi net\_manager secara keseluruhan stabil dan tidak mengalami perubahan signifikan dalam periode yang diukur.  Namun, perlu dilakukan pembanding historis untuk memastikan apakah nilai ini mencerminkan tren normal atau adanya anomali.

\begin{figure}[H]
\centering
\includegraphics[width=0.92\linewidth]{\detokenize{chart_distribution.png}}
\caption{ Distribusi frekuensi nilai pada kolom net\_manager. }
\label{fig:distribution-net-manager}
\end{figure}
\subsection{ Rata-rata per Kolom Numerik }
Analisis data energi menunjukkan rata-rata global active power sebesar 8716912000008 kW. Total energy yang dihasilkan dari sum(global\_active\_power) / 60 = 4417.023129251701 kWh.  Rata-rata jumlah koneksi terhubung mencapai 21.76582709326072, dengan delivery percentage sebesar 98.62490476190473.  Persentase koneksi aktif mencapai 99.26479591836733\%, dan persentase jenis koneksi mencapai 62.76242341729068\%. Data ini perlu dibandingkan dengan data historis untuk mengidentifikasi anomali yang mungkin terjadi.

\begin{figure}[H]
\centering
\includegraphics[width=0.92\linewidth]{\detokenize{chart_averages.png}}
\caption{ Perbandingan rata-rata antar kolom numerik. }
\label{fig:column-averages}
\end{figure}
\subsection{ Frekuensi Kategori: street }
Analisis terhadap data frekuensi kategori street menunjukkan bahwa "Oostercluft" merupakan kategori yang paling banyak dikunjungi dengan jumlah unik mencapai 704. Total jumlah pengunjung pada kategori ini mencapai 1470, menunjukkan dominasinya dalam kategori tersebut.  Untuk mendapatkan gambaran lebih lengkap mengenai dinamika penggunaan energi, diperlukan pembanding historis untuk mengukur perubahan frekuensi dan total penggunaan energi dari setiap kategori street.

\begin{figure}[H]
\centering
\includegraphics[width=0.92\linewidth]{\detokenize{chart_categories.png}}
\caption{ Distribusi 15 kategori teratas pada kolom street. }
\label{fig:categories-street}
\end{figure}

\section{Synthesis and Implications}
Secara keseluruhan, hasil visualisasi menunjukkan beberapa pola yang perlu dibaca bersama. Grafik distribusi net\_manager menunjukkan nilai yang konsisten dan seragam dengan rata-rata sebesar 8716912000008.0 kW. Analisis data energi menunjukkan rata-rata global active power sebesar 8716912000008 kW. Analisis terhadap data frekuensi kategori street menunjukkan bahwa "Oostercluft" merupakan kategori yang paling banyak dikunjungi dengan jumlah unik mencapai 704. Penilaian anomali tetap memerlukan pembanding historis atau batas operasional eksplisit.

\section{Conclusion}
Laporan ini menunjukkan bahwa pipeline lokal dapat menghasilkan grafik, statistik ringkas, dan narasi analisis secara sequential untuk dataset sembarang. Hasil ini dapat digunakan sebagai dasar eksplorasi data lanjutan, dengan catatan bahwa klaim anomali memerlukan pembanding historis tambahan.

\end{document}